\section{结果沟通与交付}

分析只有在他人能消费时才有用。根据受众选择合适的交付格式：

\begin{table}[H]
\centering
\begin{tabular}{ll}
\toprule
\textbf{交付形式} & \textbf{适用场景} \\
\midrule
Markdown memo & 技术协作者 \\
CSV / Spreadsheet & 下游运营团队 \\
\texttt{.docx} 报告 & 需要排版和表格的正式文档 \\
PDF & 最终交付物或附录 \\
静态报告站点 & 需要以 URL 分享的结果 \\
\bottomrule
\end{tabular}
\caption{不同受众对应的交付格式}
\end{table}

\begin{importantbox}{别忘了写 Caveats}
如果 join 质量不完美、存在采样偏差、或模型假设脆弱，Codex 应当在交付物中\textbf{明确说明这些局限}。诚实的 caveats 是建立信任的关键。
\end{importantbox}

\subsection{本章小结}
交付物的形式应匹配受众需求。无论选择哪种格式，都必须包含诚实的局限性说明。

